\section{Introduction}\label{sec:introduction}

The Anderson model asks how independent disorder changes the spectral
and transport properties of a lattice Hamiltonian. We prove a
weak-disorder delocalization result for independent uniform site
potentials in every fixed dimension $d\ge3$: on a fixed open interval,
the spectrum is purely absolutely continuous and has nonzero weight.
Both eigenvalue mass and singular continuous mass are excluded.

\subsection{The model and the result}

Fix an integer $d\ge3$. On the product probability space
\[
 \Omega=[-1,1]^{\ZZ^d},\qquad
 \PP=\left(\tfrac12\mathbf1_{[-1,1]}(v)\,dv\right)^{\otimes\ZZ^d},
\]
let $V_x(\omega)=\omega_x$. For $\lambda>0$, define the bounded
self-adjoint operator on $\ell^2(\ZZ^d)$ by
\begin{equation}\label{eq:model}
 (H_{\lambda,\omega}\psi)(x)
 =-\sum_{j=1}^d\bigl(\psi(x+\mathbf e_j)+\psi(x-\mathbf e_j)\bigr)
   +\lambda V_x(\omega)\psi(x).
\end{equation}
Here $\mathbf e_1,\ldots,\mathbf e_d$ are the coordinate vectors.
Multiplication by $(-1)^{x_1+\cdots+x_d}$ is a unitary operator that
reverses the hopping sign and leaves the diagonal potential unchanged.
Thus the positive-hopping convention gives the identical conclusion.
For a Borel set $B$, write $\chi_B(H)$ for the spectral projection,
and let $\mathcal H_{\mathrm{ac}}(H)$ be the subspace of vectors whose
spectral measures are absolutely continuous with respect to Lebesgue
measure.

\begin{theorem}\label{thm:main}
There is a fixed sufficiently small number $e_2\in(0,1/100)$ such that,
for every fixed integer $d\ge3$, there is a number $\lambda_d>0$
with the following property. Set
\begin{equation}\label{eq:energy-windows}
 I_d=
 \begin{cases}
   (-6+e_2/2,-6+e_2),&d=3,\\
   -2d+(1/200,1/100),&d\ge4.
 \end{cases}
\end{equation}
Then
for every fixed $0<\lambda<\lambda_d$, almost surely
\[
 I_d\subset\sigma(H_{\lambda,\omega}),\qquad
 \chi_{I_d}(H_{\lambda,\omega})\ell^2(\ZZ^d)
       \subset\mathcal H_{\mathrm{ac}}(H_{\lambda,\omega}).
\]
In particular, the spectral restriction to $I_d$ is purely absolutely
continuous and $\chi_{I_d}(H_{\lambda,\omega})\ne0$.
For $d\ge4$, one may take $\lambda_d=\sqrt3\,2^{-n_*(d)}$, where the finite-matrix
specification of $n_*(d)$ is given in Definition~\ref{def:threshold}.
For $d=3$, the proof supplies an existential positive threshold.
\end{theorem}

The interval is independent of $\lambda$. The probability-one event
may depend on the fixed coupling, and all estimates may depend on the
fixed dimension. No uniformity as $d\to\infty$ is asserted. The
integer $n_*(d)$ for $d\ge4$ is specified by an infinite family of
probability conditions on finite matrices. The proof establishes its
existence but supplies no effective numerical bound.

\subsection{Historical context and scope}

Anderson's random-site-energy model explains how spatial disorder can
suppress quantum transport~\cite{Anderson1958}. The competing
weak-disorder question is whether extended states survive on a
Euclidean lattice of dimension at least three. Simon's Problem~1
asks for a purely absolutely continuous energy range for the
nearest-neighbor model with iid uniform potentials in dimensions at
least three, at suitable disorder widths~\cite{Simon2000}.
Theorem~\ref{thm:main} supplies the purely absolutely continuous
energy range requested in that problem for every sufficiently small
fixed disorder width in each fixed dimension $d\ge3$. It does not
locate the full mobility edge discussed after Problem~1 or resolve
the diffusion question stated separately as Problem~3.

Rigorous localization was developed through several complementary
methods. Fr\"ohlich and Spencer introduced a multiscale Green-function
analysis in strong-disorder and suitable extreme-energy
regimes~\cite{FrohlichSpencer1983}; Fr\"ohlich, Martinelli, Scoppola
and Spencer subsequently established pure-point spectrum with
exponentially decaying eigenfunctions in such
regimes~\cite{FMSS1985}. The fractional-moment method of Aizenman
and Molchanov gives another route to spectral localization at large
disorder and extreme energies~\cite{AizenmanMolchanov1993}.
Corollary~\ref{cor:coexistence} combines a fixed-disorder
edge-localization theorem with Theorem~\ref{thm:main} at the same
coupling. A localized edge interval alone does not determine the
spectral type throughout the rest of the band.

On the Bethe lattice, Klein proved purely absolutely continuous
spectrum on compact intervals inside the free band at sufficiently
weak disorder~\cite{Klein1994,Klein1998}. Aizenman, Sims and Warzel
established stability of the absolutely continuous component by a
fluctuation-based method~\cite{AizenmanSimsWarzel2006}; that result
does not itself assert spectral purity. Aizenman and Warzel used
resonant delocalization to obtain an absolutely continuous component
beyond the free band for regular full-support laws at weak disorder,
and proved pure absolute continuity near the spectral edges in a weak bounded-disorder
regime that includes the uniform law~\cite{AizenmanWarzel2013}.
For sufficiently high-degree trees, at a specified hopping scale and
for a regular class of full-support single-site densities, Aggarwal
and Lopatto obtained mobility edges separating absolutely continuous
and pure-point intervals~\cite{AggarwalLopatto2025}. The branching
geometry distinguishes these results from the model on $\ZZ^d$; the
last result also has different law and scale assumptions.

Jak\v{s}i\'c, Last and Warzel gave a further tree argument and a
conditional lattice criterion through the Schur complement of a
two-site Poisson-transformed spectral measure~\cite{JaksicLastWarzel2026}.
Their criterion assumes a positive lower limit of the expected
energy-integrated quantity and concludes a nontrivial absolutely
continuous component. It does not establish that lower bound for the present
weak-disorder lattice model.

Erd\H{o}s, Salmhofer and Yau proved a weak-coupling diffusion limit for
the lattice model: at disorder-dependent long times and spatial scales,
the expected Wigner distribution converges to a heat-equation solution
as the coupling tends to zero~\cite{ErdosSalmhoferYau2007}.
For Gaussian disorder, Black, Drogin and Hern\'andez obtained
estimates at long finite times and positive resolvent heights in a
weak-coupling regime~\cite{BlackDroginHernandez2025}.
Here the boundary limit is taken at each fixed sufficiently small
positive coupling, and the conclusion also excludes singular spectral
mass.

\subsection{The proof mechanism}

The proof first establishes a finite-volume resolvent certificate.
On a large periodic lattice, we construct a set containing at least
half the sites and add small complex damping only on that set. The
diagonal imaginary parts there remain bounded below as the damping
tends to zero. At each site, the final real potential is its original
independent uniform variable. Proposition~\ref{prop:certificate}
states the certificate precisely.

We begin with the free operator and write each uniform potential as a
series of independent binary signs. Revealing a sign replaces part of
the artificial damping by true randomness. Active sites retain damping;
terminal sites have their complete true potentials and have been
eliminated by a Schur complement. During a stage, a diagonal
self-consistency parameter remains frozen until the stage-end reset.
We carry three controls for the rescaled inverse on active physical
sites. Its complex entrywise square controls the linearized equation
for the eventual diagonal reset. A bound on row fourth moments prevents
the mass of a resolvent row from concentrating at a few sites, limiting
the influence of individual updates. The Ward identity turns squared
absolute resolvent entries into a Laplacian whose inverse form supplies
a spatial energy norm for
accumulating local changes. Fourier estimates for the free resolvent
initialize these bounds with slack on the chosen energy interval.

During the one-digit stages, ordinary updates are tested on finite
tiles, since a union bound over all sites at all scales would be too
expensive. A coupled active site carries a true prefix with an
untouched suffix, and a first tile trial
uses its next true sign when that site is reached. If an ordinary tile
trial fails, its numerical changes are rolled back and every answer remains in the
retained information. All active sites in that tile are marked
\emph{pending}, including sites not reached by the trial. The mark is
conservative: at a pending site the construction ceases to rely on the
true-prefix and untouched-suffix guarantees. Ordinary updates at
pending sites use independent auxiliary digits; a retry uses new ones
throughout the tile.

At later scales we clear small components of pending sites. A unitary
scattering matrix records the interaction between a fixed core and its
exterior. The range of its cross block contains the survivor directions
through which completing the core can act. Unitarity bounds the size
of that change, so small diagonal entries of the projection onto the
cross-block range make its effect small in every surviving row. This
gives a deterministic target for preparing the core, with no freshness
assumption on the core values.

Regularized projection weights guide exposures of true residuals at
exterior sites. A spatial estimate controls the entire cross-block
range even when its nonzero singular values are arbitrarily small;
together with component separation and the retained marking rules, it
keeps selected exterior pivots close to their core and ensures that
their true residuals are still untouched. A negative drift
bounds the number of exposures, while stopped martingale estimates
control their accumulated effect, including a first violating update.
After successful preparation, any still-unknown true core values are
revealed and recorded, and the deterministic completion is applied.
The core and exposed exterior sites then become committed terminals.
A failed clearing rolls back its provisional matrix changes, keeps
every answer, and marks its core and every inspected exterior site pending.

Two inductions connect these operations to the certificate. A
deterministic induction combines committed local changes, controls the
geometry of terminal sites, keeps at least half the sites active, and
closes the self-consistent reset. A probabilistic induction tracks the
retained information and bounds failures at deterministically ordered
labels by successive conditioning. A surviving pending site yields
either a recent failure label or many separated older witnesses whose
leaf labels are distinct. With high probability the cutoff finds no
pending sites and the later
sweeps succeed. The last sweep then inserts the entire remaining true
suffix at every active site. Those sites keep the final small damping,
although their real potentials are now complete true values.

The spectral argument converts the certificate to information about
the ordinary resolvent of the true operator. A cut by its spectral
projection onto a narrow energy window controls the comparison through
that projection's rank for every qualifying certificate. Own-site
averaging under the original product law of the true potentials bounds
the expected rank, giving a positive expectation for a bounded function of the true
scalar diagonal imaginary part. That estimate passes to infinite
volume. At almost every fixed energy, ergodicity and rank-one changes
of a site potential then propagate boundary positivity to every site
almost surely. Conditioning on the operator with one site removed
produces an exceptional energy set independent of the remaining site
variable. Absolute continuity holds on its complement, and spectral
averaging removes mass on the exceptional set itself. This final step
gives purity.

The form of the construction is common to all fixed $d\ge3$.
Dimension enters through three estimates. A deliberately weaker
Fourier bound retains the three-dimensional interpolation gains in
every dimension, with a convex-slice argument handling two-resolvent
shell tangencies. The allowed size exponent for pending components
is chosen inversely proportional to $d$, keeping spatial packing
losses below the analytic gains. The history recursion uses a
dimension-dependent separation of scales to pay for the possible
positions of spatial witnesses.

Related resolvent controls appear in Black, Drogin and Hern\'andez's
work on the Gaussian Anderson model: a self-consistent diagonal shift,
concentration governed by fourth-power row sums, and an elliptic
equation for averaged squared resolvent
entries~\cite[Sections~1.1--1.2]{BlackDroginHernandez2025}.
Their spatial equation develops the $T$-equation approach of
Erd\H{o}s, Knowles, Yau and Yin for random band
matrices~\cite{ErdosKnowlesYauYin2013}. Here the Ward identity
gives an exact Laplacian for each current matrix state, and we transport
its inverse form and the other controls through the adaptive
revelations.

\subsection{Coexistence with edge localization}

The interval $I_d$ lies above the lower free band edge $-2d$.
At the lower edge of the almost-sure spectrum, fixed-disorder localization
gives a separate interval. Combining that result with
Theorem~\ref{thm:main} yields the following consequence at the same
fixed coupling.

\begin{corollary}[Spectral coexistence]\label{cor:coexistence}
For every fixed $d\ge3$ and every fixed $0<\lambda<\lambda_d$,
there is a deterministic number $\delta_{d,\lambda}\in(0,\lambda/2]$
such that the open interval
\[
 J_{d,\lambda}
 =(-2d-\lambda,-2d-\lambda+\delta_{d,\lambda})
\]
is disjoint from $I_d$ and, almost surely, both
$\chi_{I_d}(H_{\lambda,\omega})$ and
$\chi_{J_{d,\lambda}}(H_{\lambda,\omega})$ are nonzero.
The restriction to $I_d$ is purely absolutely continuous, whereas the
restriction to $J_{d,\lambda}$ is pure point with exponentially decaying
eigenfunctions. The probability-one event may depend on $d$ and $\lambda$.
\end{corollary}

\begin{proof}
The variables $\lambda V_x$ have the iid law
$\mu_\lambda(dv)=(2\lambda)^{-1}\mathbf1_{[-\lambda,\lambda]}(v)\,dv$.
Its concentration function is
$\sup_a\mu_\lambda([a,a+t])=\min\{1,t/(2\lambda)\}$, so it satisfies
the bounded-support H\"older hypothesis of the fixed-disorder edge
localization theorem of Elgart and Klein. Their author-preprint
Theorem~2.2 and Corollary~1.8(i)--(ii)~\cite{ElgartKlein2019}
give a deterministic $A>0$ and almost-sure pure-point spectrum with
exponentially decaying eigenfunctions on
$[-2d-\lambda,-2d-\lambda+A)$. Take
$\delta_{d,\lambda}=\min\{A/2,\lambda/2\}$.
The almost-sure spectrum is $[-2d-\lambda,2d+\lambda]$
\cite[author-preprint equation~(1.17)]{ElgartKlein2019}, hence its
projection on the nonempty open interval $J_{d,\lambda}$ cannot vanish.
This interval lies below $-2d$, while $I_d$ lies above $-2d$.
Intersect the localization and spectrum events with the probability-one
event of Theorem~\ref{thm:main}, at the same fixed coupling.
\end{proof}

\paragraph{Organization.}
Section~\ref{sec:state} sets up the certificate and matrix states.
Section~\ref{sec:initialization} initializes the invariant bounds.
Sections~\ref{sec:sweep} and \ref{sec:local-sweeps} establish ordinary
updates and their stopped local trials. Sections~\ref{sec:clearing}
and \ref{sec:assembly} prepare pending cores, combine the committed
changes, and close the deterministic induction.
Section~\ref{sec:history} controls the retained failures and completes
the finite-volume certificate and the threshold specification.
Section~\ref{sec:spectral} converts the certificate to
Theorem~\ref{thm:main}.
